\guidechapter{Why moving data is part of the mathematics}{Read before the machine-model section.}

Imagine a long row of labeled boxes, with a head that can move one box per
step. The label is a conceptual address; the machine does not get free random
access just because we can write $A[i,j]$ on paper. A multidimensional array
is flattened into one ordered stream. In lexicographic order, the last
coordinate changes fastest. For a $2\times3$ array the order is
\[
 (0,0),(0,1),(0,2),(1,0),(1,1),(1,2).
\]
Changing to the opposite coordinate order rearranges the records. On this
machine, the cost includes getting them into those positions.

\paragraph{A role is a place in a circuit; a row is a block of data.}
The finite network has $W$ roles, counting both data roles and scratch roles.
At a recursion node, complete rows are distributed cyclically among them.
A row contains the entire suffix of active address fields. That completeness
ensures that each role sees the same local address set. After suitable padding,
each role receives exactly a $1/W$ fraction of the logical volume.

\explfigure{02-role-volume}{The volume division that drives the recurrence.}{
Each illustrative row has $B$ bits. Twelve rows give $V=12B$; three role
streams each hold $4B$. The actual paper uses much larger fixed networks.}

\paragraph{Why auxiliary roles are included.}
Suppose we counted only useful output rows and forgot workspace. A construction
could appear cheap merely by hiding a large amount of temporary data. The
network's $W$ counts all roles. Scratch roles can hold arbitrary existing data
because the scalar circuit restores them. They are not presumed free zeroed
memory. The later framed array operation must also be correct on these roles.

\paragraph{Logical volume is not permission to ignore copying.}
$V$ names the data volume used to normalize a recursive call. Physical copies,
padding, counter updates, tape-head returns, and cleanup still consume time.
They are charged as overhead. When a child is active, the paper's schedule
parks its parent's streams and restores them on return. The claimed accounting
requires that it does not repeatedly sweep through every parked ancestor.

\paragraph{A fixed machine can still process growing descriptors.}
Input lengths, axis lengths, and call states are stored as data. The number
of tapes and the alphabet cannot grow with $n$. This is why a proof about a
finite circuit is only one ingredient: it must be transferred to uniform
stream operations with explicit costs for manipulating those descriptors.

\takeaway{Watch both the number of recursive calls and the amount of data
each call processes. A saving in one can disappear if the other is miscounted.}
\sourceanchor{02-streams.tex, lem:elementary-streams}
